
\section{What causes the \nn{}}
\label{sec:why}

There are two main areas that I believe allows this to happen.
The developers and the controller-runtime together allow the issue to araise.
Not one item under the areas is enough to causes the \nn{}, but together they both cause a perfect storm allowing the \nn{} to appear.

\subsection{Developers and their development environment}
The developer environment I believe is one of the main reason the \nn{} is not spotted until it is too late, when our operators are in the hands of our customers.
I know from my own set-up, my team tries to be resourcefully with resources, and enable development to be carried out on local clusters such as KinD\cite{noauthor_kind_nodate}.
This allows us to build clean development environments quickly, and cheaply.
We do not reuse existing clusters.
This reduces the development costs as long-running clusters, and the needed maintenance is not required.

The \nn{} only happens when there is other third party resources on the cluster, such as customer workloads.
Even then, a numerous number of the third party resources is required.
As the development cluster are constantly recreated the increased memory usage from the third party resource are never detected.
Even our release testing cluster tend to be freshly created for that test cycle, and don't have the legacy resources a customers cluster would have.

\subsection{controller-runtime, and it's defaults}
The controller-runtime is the default package that most operators written in golang\cite{noauthor_go_nodate} will use to communicate with the API server.
It is very good at getting an operator bootstrapped, but has a number of default shortcoming's that are misleading developers.

The first behaviour that is misleading is what the predicate functions do for the reconcile functions in the reconcile loop.
Or more so the assumption developers make about these predicate functions.
Most of the developers I have asked about the behaviour all assumed that if an event failed to pass the predicate function, the resources related to that event are not cached.
While in truth the caching of the resources happens before the predicate function is ever triggered.

The next behaviour that is unexpected is related to the watches the controller-runtime configures.
When defining a reconcile function the developer will also define what resources will trigger that reconcile function.
The predicate functions are also defined at this time.
Speaking with other developers there is an understanding that the controller-runtime will configure watches on the defined resources of the reconcile function.
It does, and this is acceptable behaviour.
The behaviour that is not expected, and not known by most I spoke with is when the standard client configured thought the controller-runtime is used to get a resource from the cluster, the controller-runtime also sets up a watch on that resources.
Even if the resource is of a known name and namespace, the controller-runtime configures the watch.
The controller-runtime can not know that the one requested is only ever accessed once.

The last behaviour that I believe leads to operator being misconfigured is lack of requiring the user to configure the cache for the controller-runtime.
Awareness of the cache configuration is missing by developers.
If it is not already configured, the developers forget to configure it.
The lack of forcing the configuration of the cache turns the cache configuration into a second class citizen.

\phantomsection\label{sec:why:controller-runtime:startup}
But the controller-runtime default start up behaviour is the kicker that makes the \textit{"why the \nn{} is such a problem"} (see \cref{sec:why:problem}) even more of a problem.
On the operators start up, the controller-runtime waits for all the events for the resources that watches have being configured for before triggering any of the reconcile functions.
This is a great behaviour as it allows the reconcile functions to act quickly by not needing to make network requests to the cluster.
In the case of the \nn{} this means before the operator can do any work all state most be known.
